\chapter{Elevated Liver Enzymes}

\section*{Definition}
Elevated liver enzymes are blood-test abnormalities, most commonly increased alanine aminotransferase (ALT), aspartate aminotransferase (AST), alkaline phosphatase (ALP), or gamma-glutamyl transferase (GGT). They indicate possible liver or biliary injury but do not, by themselves, measure liver function or establish a diagnosis.

\section*{What Happens in Your Body}
ALT and AST rise when hepatocytes are injured and release intracellular enzymes. ALP and GGT more strongly suggest cholestasis or bile-duct injury, although ALP may also originate from bone. The pattern, degree, duration, symptoms, medications, and synthetic markers such as bilirubin, albumin, and INR determine urgency.

\section*{Causes}
\begin{itemize}
  \item Metabolic dysfunction-associated steatotic liver disease, alcohol-associated liver disease, and viral hepatitis.
  \item Medication or supplement injury, including acetaminophen overdose, selected antibiotics, anticonvulsants, statins, herbal products, and bodybuilding supplements.
  \item Biliary obstruction from gallstones, strictures, or tumors; autoimmune hepatitis, primary biliary cholangitis, and primary sclerosing cholangitis.
  \item Ischemic or congestive injury, muscle injury causing AST elevation, hemochromatosis, Wilson disease, and alpha-1 antitrypsin deficiency.
\end{itemize}

\section*{When to See a Doctor}
Seek emergency care for confusion, severe drowsiness, vomiting blood, black stools, fainting, severe right-upper-abdominal pain with fever, or jaundice with rapidly worsening illness. Prompt assessment is needed for new jaundice, dark urine, pale stools, persistent vomiting, or rapidly increasing results.

\section*{Related Symptoms}
\begin{itemize}
  \item Fatigue, nausea, reduced appetite, right-upper-abdominal discomfort, itching, jaundice, dark urine, or pale stools.
  \item Many people have no symptoms and are identified through routine blood testing.
\end{itemize}

\section*{Self-Care / Home Management}
Avoid alcohol and unnecessary supplements. Do not exceed labeled or prescribed acetaminophen doses, and do not stop essential prescription medicines without medical advice. Bring a complete medication, supplement, alcohol, travel, and hepatitis-risk history to evaluation.

\section*{How Your Doctor Will Evaluate You}
\begin{itemize}
  \item Repeat and fractionate the panel: ALT, AST, ALP, GGT, bilirubin, albumin, and INR; review the pattern and trend.
  \item Test for viral hepatitis and, when indicated, autoimmune disease, iron overload, Wilson disease, and alpha-1 antitrypsin deficiency.
  \item Ultrasound is commonly used for steatosis, masses, gallstones, and bile-duct dilation; additional imaging or biopsy is selective.
\end{itemize}

\section*{Treatment Options}
Treatment addresses the cause: weight reduction and metabolic-risk management for steatotic liver disease, alcohol cessation, antiviral therapy for selected viral hepatitis, withdrawal of a causative drug under supervision, and urgent biliary decompression when obstruction is present. Acute liver failure requires hospital and specialist care.

\section*{References}
\begin{thebibliography}{99}
\bibitem{elevated_liver_enzymes:ref1} Kwo PY, Cohen SM, Lim JK. ACG Clinical Guideline: Evaluation of abnormal liver chemistries. \textit{Am J Gastroenterol}. 2017;112:18--35.
\bibitem{elevated_liver_enzymes:ref2} European Association for the Study of the Liver. EASL clinical practice guidelines on non-invasive liver disease assessment. \textit{J Hepatol}. 2021;75:659--689.
\end{thebibliography}
